\section{Les bases, sans prérequis}
\frame{\sectionpage}

\begin{frame}{Le langage dans le cerveau}
\begin{itemize}
    \item Un \textbf{réseau du langage} fronto-temporal, principalement dans l'hémisphère gauche
    \item Il répond à la langue, écrite ou parlée, et peu à autre chose
    \item Il est distinct du réseau de la pensée abstraite et du raisonnement
\end{itemize}
% IMAGE : cerveau avec le réseau du langage (à ajouter)
\src{Fedorenko et al. ; Fedorenko, Open Encyclopedia of Cognitive Science, 2025}
\end{frame}

\begin{frame}{Comment « écouter » le cerveau ?}
\centering\small
\begin{tabular}{lcccc}
    \textcolor{yellow}{Méthode} & \textcolor{yellow}{Espace} & \textcolor{yellow}{Temps} & \textcolor{yellow}{Invasif ?} \\
    \hline
    IRM fonctionnelle (fMRI) & millimètres & secondes & non \\
    EEG / MEG & centimètres & millisecondes & non \\
    Intracrânien (iEEG, sEEG) & millimètres--cm & millisecondes & oui \\
    Neurones uniques & micromètres & millisecondes & oui \\
\end{tabular}
\vspace{1em}

\normalsize
L'intracrânien se fait chez des patients épileptiques implantés pour raisons cliniques.
% IMAGE : schémas des méthodes (à ajouter)
\end{frame}

\begin{frame}{Le protocole type}
\begin{enumerate}
    \item Un participant lit ou écoute un texte (souvent une histoire naturelle)
    \item On enregistre son activité cérébrale mot par mot
    \item On donne le même texte à un modèle de langage
    \item On compare les deux
\end{enumerate}
\end{frame}

\begin{frame}{Comparer un modèle à un cerveau}
\centering
\begin{tikzpicture}[
    box/.style={draw=yellow, rounded corners, fill=darkred, text=white, align=center, minimum height=1.1cm, font=\small},
    arr/.style={-{Stealth}, thick, yellow}]
    \node[box] (txt) at (0,0) {Texte\\{\footnotesize « Le chat dort… »}};
    \node[box] (lm) at (4.6,1.1) {Modèle de langage\\{\footnotesize activations}};
    \node[box] (br) at (4.6,-1.1) {Cerveau\\{\footnotesize fMRI / EEG / iEEG}};
    \node[box] (reg) at (9.2,0) {Régression\\linéaire};
    \node[box] (sc) at (12.2,0) {Score\\d'alignement};
    \draw[arr] (txt) -- (lm);
    \draw[arr] (txt) -- (br);
    \draw[arr] (lm) -- (reg);
    \draw[arr] (br) -- (reg);
    \draw[arr] (reg) -- (sc);
\end{tikzpicture}

\vspace{1em}
\begin{itemize}
    \item On prédit le signal cérébral à partir des activations du modèle
    \item Le score est mesuré sur des données \textbf{non vues} pendant l'ajustement
    \item On le compare à un \textbf{plafond de bruit} (la variabilité entre participants)
\end{itemize}
\end{frame}

\begin{frame}{Deux outils à retenir}
\begin{block}{Modèle d'encodage}
    Régression des activations du modèle vers le signal cérébral (utilisé dans Brain-Score).
\end{block}
\begin{block}{Analyse de similarité représentationnelle (RSA)}
    On compare les distances entre stimuli dans le modèle à celles mesurées dans le cerveau.
\end{block}
\end{frame}
